\secciongrande{Programación de scripts dónde y cómo colocar los scripts}

\resumen{
  Los scripts de javascript se insertan en un documento html utilizando la etiqueta \texttt{<script>}. Pueden ubicarse tanto en el encabezado (entre \texttt{<head>} y \texttt{</head>}) como en el cuerpo (entre \texttt{<body>} y \texttt{</body>}). La etiqueta de apertura debe incluir el atributo \texttt{type="text/javascript"} para indicar el formato del código.
}

\definicion{Etiqueta script}{Etiqueta contenedora html que designa el inicio y el fin de un bloque de código javascript, comenzando con \texttt{<script>} y terminando con \texttt{</script>}.}
\subtituloA{Ubicación de los scripts en html}

\lista{
  \textbf{Encabezado:} Se coloca entre \texttt{<head>} y \texttt{</head>}. |
  \textbf{Cuerpo:} Se coloca entre \texttt{<body>} y \texttt{</body>}.
}
